\begin{pastexam}[label={px:ncku-109-prog-07}]{成大資訊聯招 109 學年度 演算法第 4 題}
（計算）

Determine an Longest Common Subsequence of <1,0,0,1,0,1,0,1> and <0,1,0,1,1,0,1,1,0>. (10\%)
\end{pastexam}

\begin{pastsoln}
答案：LCS 的長度是 $6$，其中一個是 $\langle 1, 0, 0, 1, 1, 0\rangle$。

令 $X = \langle 1,0,0,1,0,1,0,1\rangle$（$m = 8$）、$Y = \langle 0,1,0,1,1,0,1,1,0\rangle$
（$n = 9$）。 $c[i, j]$ 是 $X$ 的前 $i$ 個和 $Y$ 的前 $j$ 個的 LCS 長度，照資結書第~18.3~節的
遞迴式填表：
\[
c[i, j] =
\begin{cases}
0 & i = 0 \text{ 或 } j = 0,\\
c[i-1, j-1] + 1 & i, j > 0 \text{ 且 } x_i = y_j,\\
\max(c[i-1, j],\, c[i, j-1]) & i, j > 0 \text{ 且 } x_i \ne y_j.
\end{cases}
\]

一列一列、每列由左到右填（資結書虛擬碼~18.6）。 粗體是回溯時經過的「相等」格：
\begin{center}
\small
\begin{tabular}{cc|cccccccccc}
& $j$ & 0 & 1 & 2 & 3 & 4 & 5 & 6 & 7 & 8 & 9 \\
$i$ & $x_i \backslash y_j$ & & 0 & 1 & 0 & 1 & 1 & 0 & 1 & 1 & 0 \\ \hline
0 & & 0 & 0 & 0 & 0 & 0 & 0 & 0 & 0 & 0 & 0 \\
1 & 1 & 0 & 0 & \textbf{1} & 1 & 1 & 1 & 1 & 1 & 1 & 1 \\
2 & 0 & 0 & 1 & 1 & \textbf{2} & 2 & 2 & 2 & 2 & 2 & 2 \\
3 & 0 & 0 & 1 & 1 & 2 & 2 & 2 & \textbf{3} & 3 & 3 & 3 \\
4 & 1 & 0 & 1 & 2 & 2 & 3 & 3 & 3 & \textbf{4} & 4 & 4 \\
5 & 0 & 0 & 1 & 2 & 3 & 3 & 3 & 4 & 4 & 4 & 5 \\
6 & 1 & 0 & 1 & 2 & 3 & 4 & 4 & 4 & 5 & \textbf{5} & 5 \\
7 & 0 & 0 & 1 & 2 & 3 & 4 & 4 & 5 & 5 & 5 & \textbf{6} \\
8 & 1 & 0 & 1 & 2 & 3 & 4 & 5 & 5 & 6 & 6 & 6 \\
\end{tabular}
\end{center}
舉兩格說明怎麼填：$c[4, 7]$，$x_4 = 1 = y_7$，取左上 $c[3, 6] + 1 = 4$。 $c[5, 7]$，
$x_5 = 0 \ne y_7 = 1$，取上面 $c[4, 7] = 4$ 和左邊 $c[5, 6] = 4$ 的較大者，$4$。

右下角 $c[8, 9] = 6$，LCS 長度是 $6$。

\textbf{回溯}（資結書虛擬碼~18.7）：從 $(8, 9)$ 出發，$x_i = y_j$ 就記下這個字元、往左上走；
不相等時，往 $c$ 值較大的方向走，上面和左邊一樣大時往上走。
\begin{itemize}
\item $(8,9)$：$x_8 = 1 \ne y_9 = 0$，上 $c[7,9] = 6 \ge$ 左 $c[8,8] = 6$，往上。
\item $(7,9)$：$x_7 = 0 = y_9$，記下 $0$，到 $(6,8)$。
\item $(6,8)$：$x_6 = 1 = y_8$，記下 $1$，到 $(5,7)$。
\item $(5,7)$：$x_5 = 0 \ne y_7 = 1$，上 $c[4,7] = 4 \ge$ 左 $c[5,6] = 4$，往上。
\item $(4,7)$：$x_4 = 1 = y_7$，記下 $1$，到 $(3,6)$。
\item $(3,6)$：$x_3 = 0 = y_6$，記下 $0$，到 $(2,5)$。
\item $(2,5)$：$x_2 = 0 \ne y_5 = 1$，上 $c[1,5] = 1 <$ 左 $c[2,4] = 2$，往左；$(2,4)$ 同理往左。
\item $(2,3)$：$x_2 = 0 = y_3$，記下 $0$，到 $(1,2)$。
\item $(1,2)$：$x_1 = 1 = y_2$，記下 $1$，到 $(0,1)$，結束。
\end{itemize}
記下的字元倒過來是 $\langle 1, 0, 0, 1, 1, 0\rangle$。 核對：它在 $X$ 裡是第 $1, 2, 3, 4, 6, 7$
個，在 $Y$ 裡是第 $2, 3, 6, 7, 8, 9$ 個，兩邊的順序都保持。

題目說「an」LCS，寫出一個就可以。 長度 $6$ 的公共子序列不只一個，例如
$\langle 0,1,0,1,0,1\rangle$、$\langle 1,0,1,0,1,0\rangle$ 也是（共 $8$ 個不同的），回溯時遇到相等
往左走就會得到別的答案。 不管寫哪一個，長度都必須是 $6$。
\end{pastsoln}
